\begin{textsample}{POS Dim 1 – 2001 – Score 1.00 – 2001/t000092.txt}  \label{ex:f1_pos_077}
I coined my own definition of success in 1934, when I was teaching at a high school in South Bend, Indiana, being a little bit disappointed, and [ disillusioned ] perhaps, by \textbf{the} way parents of \textbf{the} youngsters in my English classes expected their youngsters to get an A or a B.
They thought a C was all right for \textbf{the} neighbors ' children, because they were all average.
But they weren ' t satisfied when their own — it would make \textbf{the} teacher feel that they had failed, or \textbf{the} youngster had failed.
And that ' s not right. \textbf{The} good Lord in his infinite wisdom didn ' t create us all equal as far as intelligence is concerned, any more than we ' re equal for size, appearance.
Not everybody could earn an A or a B, and I didn ' t like that way of judging, and I did know how \textbf{the} alumni of various schools back in \textbf{the} ' 30s judged coaches and athletic teams.
If you won them all, you were considered to be reasonably successful — not completely.
Because I found out — we had a number of years at UCLA where we didn ' t lose a game.
But it seemed that we didn ' t win each individual game by \textbf{the} margin that some of our alumni had predicted — And quite frequently I really felt that they had backed up their predictions in a more materialistic manner.
But that was true back in \textbf{the} 30s, so I understood that.
But I didn ' t like it, I didn ' t agree with it.
I wanted to come up with something I hoped could make me a better teacher, and give \textbf{the} youngsters under my supervision, be it in athletics or \textbf{the} English classroom, something to which to aspire, other than just a higher mark in \textbf{the} classroom, or more points in some athletic contest.
I thought about that for quite a spell, and I wanted to come up with my own definition.
I thought that might help.
And I knew how Mr.
Webster defined it, as \textbf{the} accumulation of material possessions or \textbf{the} attainment of a position of power or prestige, or something of that sort, worthy accomplishments perhaps, but in my opinion, not necessarily indicative of success.
So I wanted to come up with something of my own.
And I recalled — I was raised on a small farm in Southern Indiana, and Dad tried to teach me and my brothers that you should never try to be better than someone else.
I ' m sure at \textbf{the} time he did that, I didn ' t — it didn ' t — well, somewhere, I guess in \textbf{the} hidden recesses of \textbf{the} mind, it popped out years later.
Never try to be better than someone else, always learn from others.
Never cease trying to be \textbf{the} best you can be — that ' s under your control.
If you get too engrossed and involved and concerned in regard to \textbf{the} things over which you have no control, it will adversely affect \textbf{the} things over which you have control.
Then I ran across this simple verse that said,"At God ' s footstool to confess, a poor soul knelt, and bowed his head. ' I failed! ' he cried. \textbf{The} Master said, ' Thou didst thy best, that is success. '"From those things, and one other perhaps, I coined my own definition of success, which is : Peace of mind attained only through self-satisfaction in knowing you made \textbf{the} effort to do \textbf{the} best of which you ' re capable.
I believe that ' s true.
If you make \textbf{the} effort to do \textbf{the} best of which you ' re capable, trying to improve \textbf{the} situation that exists for you, I think that ' s success, and I don ' t think others can judge that ; it ' s like character and reputation — your reputation is what you ' re perceived to be ; your character is what you really are.
And I think that character is much more important than what you are perceived to be.
You ' d hope they ' d both be good, but they won ' t necessarily be \textbf{the} same.
Well, that was my idea that I was going to try to get across to \textbf{the} youngsters.
I ran across other things.
I love to teach, and it was mentioned by \textbf{the} previous speaker that I \textbf{enjoy} poetry, and I dabble in it a bit, and love it.
There are some things that helped me, I think, be better than I would have been.
I know I ' m not what I ought to be, what I should be, but I think I ' m better than I would have been if I hadn ' t run across certain things.
One was just a little verse that said,"No written word, no spoken plea can teach our youth what they should be ; nor all \textbf{the} books on all \textbf{the} shelves — it ' s what \textbf{the} teachers are themselves."That made an impression on me in \textbf{the} 1930s.
And I tried to use that more or less in my teaching, whether it be in sports, or whether it be in \textbf{the} English classroom.
I love poetry and always had an interest in that somehow.
Maybe it ' s because Dad used to read to us at night, by \textbf{coal} \textbf{oil} \textbf{lamp} — we didn ' t have electricity in our farm home.
And Dad would read poetry to us.
So I always liked it.
And about \textbf{the} same time I ran across this one verse, I ran across another one.
Someone asked a lady teacher why she taught, and after some time, she said she wanted to think about that.
Then she came up and said,"They ask me why I teach, and I reply, ' Where could I find such splendid company? ' There sits a statesman, strong, unbiased, wise ; another Daniel Webster, silver-tongued.
A doctor sits beside him, whose quick, steady hand may mend a bone, or stem \textbf{the} life-blood ' s flow.
And there a builder ; upward rise \textbf{the} arch of a church he builds, wherein that minister may speak \textbf{the} word of God, and lead a stumbling soul to touch \textbf{the} Christ.
And all about, a gathering of teachers, farmers, merchants, laborers — those who work and vote and build and plan and pray into a great tomorrow.
And I may say, I may not see \textbf{the} church, or hear \textbf{the} word, or eat \textbf{the} food their hands may grow, but yet again I may ; And later I may say, I knew him once, and he was weak, or strong, or bold or proud or gay.
I knew him once, but then he was a boy.
They ask me why I teach and I reply, ' Where could I find such splendid company? '"And I believe \textbf{the} teaching profession — it ' s true, you have so many youngsters, and I ' ve got to think of my youngsters at UCLA — 30-some attorneys, 11 dentists and doctors, many, many teachers and other professions.
And that gives you a great deal of pleasure, to see them go on.
I always tried to make \textbf{the} youngsters feel that they ' re there to get an education, number one ; basketball was second, because it was paying their way, and they do need a little time for social activities, but you let social activities take a little precedence over \textbf{the} other two, and you ' re not going to have any very long.
So that was \textbf{the} idea that I tried to get across to \textbf{the} youngsters under my supervision.
I had three rules, pretty much, that I stuck with practically all \textbf{the} time.
I ' d learned these prior to coming to UCLA, and I decided they were very important.
One was"Never be late."Later on I said certain things — \textbf{the} players, if we were leaving for somewhere, had to be neat and clean.
There was a time when I made them wear jackets and shirts and ties.
Then I saw our chancellor coming to school in denims and turtlenecks, and thought, it ' s not right for me to keep this other [ rule ] so I let them just — they had to be neat and clean.
I had one of my greatest players that you probably heard of, Bill Walton.
He came to catch \textbf{the} bus ; we were leaving for somewhere to play.
And he wasn ' t clean and neat, so I wouldn ' t let him go.
He couldn ' t get on \textbf{the} bus, he had to go home and get cleaned up to get to \textbf{the} airport.
So I was a stickler for that.
I believed in that.
I believe in time ; very important.
I believe you should be on time, but I felt at practice, for example — we start on time, we close on time. \textbf{The} youngsters didn ' t have to feel that we were going to keep them over.
When I speak at coaching clinics, I often tell young coaches — and at coaching clinics, more or less, they ' ll be \textbf{the} younger coaches getting in \textbf{the} profession.
Most of them are young, you know, and probably newly-married.
And I tell them,"Don ' t run practices late, because you ' ll go home in a bad mood, and that ' s not good, for a young married man to go home in a bad mood.
When you get older, it doesn ' t make any difference, but —"So I did believe : on time.
I believe starting on time, and I believe closing on time.
And another one I had was, not one word of profanity.
One word of profanity, and you are out of here for \textbf{the} day.
If I see it in a game, you ' re going to come out and sit on \textbf{the} bench.
And \textbf{the} third one was, never criticize a teammate.
I didn ' t want that.
I used to tell them I was paid to do that.
That ' s my job.
I ' m paid to do it.
Pitifully poor, but I am paid to do it.
Not like \textbf{the} coaches today, for gracious sakes, no.
It ' s a little different than it was in my day.
Those were three things that I stuck with pretty closely all \textbf{the} time.
And those actually came from my dad.
That ' s what he tried to teach me and my brothers at one time.
I came up with a pyramid eventually, that I don ' t have \textbf{the} time to go on that.
But that helped me, I think, become a better teacher.
It ' s something like this : And I had blocks in \textbf{the} pyramid, and \textbf{the} cornerstones being industriousness and enthusiasm, working hard and \textbf{enjoying} what you ' re doing, coming up to \textbf{the} apex, according to my definition of success.
And right at \textbf{the} top, faith and patience.
And I say to you, in whatever you ' re doing, you must be patient.
You have to have patience to — we want things to happen.
We talk about our youth being impatient a lot, and they are.
They want to change everything.
They think all change is progress.
And we get a little older — we sort of let things go.
And we forget there is no progress without change.
So you must have patience, and I believe that we must have faith.
I believe that we must believe, truly believe.
Not just give it word service, believe that things will work out as they should, providing we do what we should.
I think our tendency is to hope things will turn out \textbf{the} way we want them to much of \textbf{the} time, but we don ' t do \textbf{the} things that are necessary to make those things become reality.
I worked on this for some 14 years, and I think it helped me become a better teacher.
But it all revolved around that original definition of success.
You know, a number of years ago, there was a Major League Baseball umpire by \textbf{the} name of George Moriarty.
He spelled Moriarty with only one ' i '.
I ' d never seen that before, but he did.
Big league baseball players — they ' re very perceptive about those things, and they noticed he had only one ' i ' in his name.
You ' d be surprised how many also told him that that was one more than he had in his head at various times.
But he wrote something where I think he did what I tried to do in this pyramid.
He called it"\textbf{The} Road Ahead, or \textbf{the} Road Behind."He said,"Sometimes I think \textbf{the} Fates must grin as we denounce them and insist \textbf{the} only reason we can ' t win, is \textbf{the} Fates themselves have missed.
Yet there lives on \textbf{the} ancient claim : we win or lose within ourselves. \textbf{The} shining trophies on our shelves can never win tomorrow ' s game.
You and I know deeper down, there ' s always a chance to win \textbf{the} crown.
But when we fail to give our best, we simply haven ' t met \textbf{the} test, of giving all and saving none until \textbf{the} game is really won ; of showing what is meant by grit ; of playing through when others quit ; of playing through, not letting up.
It ' s bearing down that wins \textbf{the} cup.
Of dreaming there ' s a goal ahead ; of hoping when our dreams are dead ; of praying when our hopes have fled ; yet losing, not afraid to fall, if, bravely, we have given all.
For who can ask more of a man than giving all within his span.
Giving all, it seems to me, is not so far from victory.
And so \textbf{the} Fates are seldom wrong, no matter how they twist and wind.
It ' s you and I who make our fates — we open up or close \textbf{the} gates on \textbf{the} road ahead or \textbf{the} road behind."Reminds me of another set of threes that my dad tried to get across to us : Don ' t whine.
Don ' t complain.
Don ' t make excuses.
Just get out there, and whatever you ' re doing, do it to \textbf{the} best of your ability.
And no one can do more than that.
I tried to get across, too, that — my opponents will tell you — you never heard me mention winning.
Never mention winning.
My idea is that you can lose when you outscore somebody in a game, and you can win when you ' re outscored.
I ' ve felt that way on certain occasions, at various times.
And I just wanted them to be able to hold their head up after a game.
I used to say that when a game is over, and you see somebody that didn ' t know \textbf{the} outcome, I hope they couldn ' t tell by your actions whether you outscored an opponent or \textbf{the} opponent outscored you.
That ' s what really matters : if you make an effort to do \textbf{the} best you can regularly, \textbf{the} results will be about what they should be.
Not necessarily what you ' d want them to be but they ' ll be about what they should ; only you will know whether you can do that.
And that ' s what I wanted from them more than anything else.
And as time went by, and I learned more about other things, I think it worked a little better, as far as \textbf{the} results.
But I wanted \textbf{the} score of a game to be \textbf{the} byproduct of these other things, and not \textbf{the} end itself.
I believe it was one great philosopher who said — no, no — Cervantes.
Cervantes said,"\textbf{The} journey is better than \textbf{the} end."And I like that.
I think that it is — it ' s getting there.
Sometimes when you get there, there ' s almost a let down.
But it ' s \textbf{the} getting there that ' s \textbf{the} fun.
As a basketball coach at UCLA, I liked our practices to be \textbf{the} journey, and \textbf{the} game would be \textbf{the} end, \textbf{the} end result.
I liked to go up and sit in \textbf{the} stands and watch \textbf{the} players play, and see whether I ' d done a decent job during \textbf{the} week.
There again, it ' s getting \textbf{the} players to get that self-satisfaction, in knowing that they ' d made \textbf{the} effort to do \textbf{the} best of which they are capable.
Sometimes I ' m asked who was \textbf{the} best player I had, or \textbf{the} best teams.
I can never answer that.
As far as \textbf{the} individuals are concerned — I was asked one time about that, and they said,"Suppose that you, in some way, could make \textbf{the} perfect player.
What would you want?"And I said,"Well, I ' d want one that knew why he was at UCLA : to get an education, he was a good student, really knew why he was there in \textbf{the} first place.
But I ' d want one that could play, too.
I ' d want one to realize that defense usually wins championships, and who would work hard on defense.
But I ' d want one who would play offense, too.
I ' d want him to be unselfish, and look for \textbf{the} pass first and not shoot all \textbf{the} time.
And I ' d want one that could pass and would pass.
I ' ve had some that could and wouldn ' t, and I ' ve had some that would and could.
So, yeah, I ' d want that.
And I wanted them to be able to shoot from \textbf{the} outside.
I wanted them to be good inside too.
I ' d want them to be able to rebound well at both ends, too.
Why not just take someone like Keith Wilkes and let it go at that.
He had \textbf{the} qualifications.
Not \textbf{the} only one, but he was one that I used in that particular category, because I think he made \textbf{the} effort to become \textbf{the} best.
There was a couple.
I mention in my book,"They Call Me Coach,"two players that gave me great satisfaction, that came as close as I think anyone I ever had to reach their full potential : one was Conrad Burke, and one was Doug McIntosh.
When I saw them as freshmen, on our freshmen team — freshmen couldn ' t play varsity when I taught.
I thought,"Oh gracious, if these two players, either one of them"— they were different years, but I thought about each one at \textbf{the} time he was there —"Oh, if he ever makes \textbf{the} varsity, our varsity must be pretty miserable, if he ' s good enough to make it."And you know, one of them was a starting player for a season and a half. \textbf{The} other one, his next year, played 32 minutes in a national championship game, Did a tremendous job for us. \textbf{The} next year, he was a starting player on \textbf{the} national championship team, and here I thought he ' d never play a minute, when he was — so those are \textbf{the} things that give you great joy, and great satisfaction to see.
Neither one of those youngsters could shoot very well.
But they had outstanding shooting percentages, because they didn ' t force it.
And neither one could jump very well, but they kept good position, and so they did well rebounding.
They remembered that every shot that ' s taken, they assumed would be missed.
I ' ve had too many stand around and wait to see if it ' s missed, then they go and it ' s too late, somebody else is in there ahead of them.
They weren ' t very quick, but they played good position, kept in good balance.
And so they played pretty good defense for us.
So they had qualities that — they came close to — as close to reaching possibly their full potential as any players I ever had.
So I consider them to be as successful as Lewis Alcindor or Bill Walton, or many of \textbf{the} others that we had ; there were some outstanding players.
Have I rambled enough?
I was told that when he makes his appearance, I was supposed to shut up.

% matched lemmas: coal, enjoy, lamp, oil, the
\end{textsample}
